% =====================================================================
% Introduction  (~320 words)
% =====================================================================
\section{Introduction}
\label{sec:introduction}

Industrial robots have become a fixture of modern manufacturing.
Worldwide annual shipments of industrial manipulators have exceeded
half a million units every year since 2021, with a global operational
stock above 3.5\,million~\cite{ifr2024}. Articulated arms with three
to seven revolute joints dominate the installed base because their
workspace flexibility is well-suited to pick-and-place, assembly and
machine tending~\cite{siciliano2009}. Increasingly these tasks are
operated remotely or autonomously, particularly in hazardous,
hygienic, or space-constrained settings where co-locating an operator
with the manipulator is undesirable.

The coursework brief~\cite{ee6003brief} reflects this trend by
requiring a two-unit embedded architecture: a handheld controller
(transmitter) that maps operator input to digital command packets, and
a robotic arm (receiver) of at least three DOF plus gripper that
executes those packets in real time and additionally supports a
fully autonomous \emph{Sense--Think--Act} sort task. The brief
mandates closed-loop joint feedback, a wireless link, mode switching
between manual and autonomous control, and explicit
performance evaluation against accuracy, latency, repeatability and
robustness metrics.

This report documents the group's design and implementation against
those requirements. Section~\ref{sec:literature} surveys the relevant
literature on industrial manipulators, low-power wireless protocols
and closed-loop control. Section~\ref{sec:design} presents the system
architecture and justifies the principal hardware choices.
Sections~\ref{sec:hardware} and~\ref{sec:software} document the
hardware and firmware respectively, including the custom packet
protocol, the per-joint PID, the autonomous-mode finite state machine
(FSM), and the computer-vision subsystem. Section~\ref{sec:cad}
covers the mechanical design and 3D printing.
Section~\ref{sec:testing} describes the test methodology and
results. Section~\ref{sec:teamwork} captures team organisation and
individual contributions, and Section~\ref{sec:conclusion} concludes.

In line with the brief, every technical section explicitly identifies
which work satisfies a \emph{core} requirement (capped at 60\,\%) and
which constitutes an \emph{advanced} embedded-systems feature
contributing to marks above that ceiling.
